%=============================================================================
% DISCUSSION AND CONCLUSION
%=============================================================================
\section{Discussion}
\label{sec:discussion}

\textbf{Where physics helps a graph network.} On this target, physics cannot
be the forecaster. A pure SEIR decoder inherits the unpredictability of its
force of infection, and 14--16\% of targets lie out of its reach. It does help
as \emph{structure}: gated onto persistence, the SEIR head makes every encoder
we tried usable and interchangeable, repairing three published architectures
that fail on their own. Section~\ref{sec:rescue} measures how much of that is the
simulator and how much the anchor it is gated onto. Coupling districts inside
the dynamics, as metapopulation GNNs do, does not help here, because
transmission has almost no local spatial structure at district-week resolution.
These results agree with our earlier work on the benchmark array, which found
that mechanistic knowledge pays when it restricts a model, not when it supplies
the prediction.

\textbf{Why the ceiling is not architectural.} The target's own last value
explains $r^2=0.85$ of the next, the best covariate 0.02, and what remains is a
move the inputs do not reveal. Every model family converges to the same
function, with residual correlation 0.87--0.99. Architecture, embeddings,
climate, synthetic data and physics penalties each change validation by less
than the seed spread. A study that searches this space and reports its best
cell will find a winner; a paired study finds noise. The NB likelihood is the
exception, because it corrects what the objective targets rather than what the
model sees.

\textbf{What the evaluation taught us.} The benchmark array leaks future data
through its row order and its climate channels. Our own first SEIR-GNN pipeline
reported a win that did not survive a correct training loop. A validation span
before each origin selects models that underperform anchored ones on test, in
five grids. Two nine-origin confirmations disagree with each other. Each of these
would have produced a publishable positive result had it gone unchecked.

\textbf{Limitations.} One country and one disease at district-week resolution.
Three origins cannot support claims at the origin unit, and our nine-origin
test spans were used twice, so they are not fresh. The two confirmations named
different endpoints. Some earlier diagnostics (renewal predictability, Moran's
$I$, the STGAT spatial penalty) were measured on the benchmark array and have
not been rerun on the corrected data. The graph is the benchmark's hand-built
adjacency; mobility data might couple districts differently.

\textbf{Next steps.} Clean confirmation needs data nothing here has touched: the
weekly reports after March 2024. Before any such run we would pre-register a
single endpoint, selection over several validation spans rather than one, and
the gated SEIR-GNN, $B$ and SEIR-LSTM as the only finalists.

\section{Conclusion}
We asked whether replacing the LSTM of an SEIR-informed forecaster with a graph
network improves district-level dengue forecasts. Under a leak-free protocol on
a corrected dataset, the answer is narrower than either literature suggests. The
graph adds little on its own (0.07). The graph's advantage over the LSTM inside
the SEIR head is consistent on three origins and not robust on nine. The largest
controlled gain comes from the likelihood, not the architecture. Physics fails
as a decoder for measurable, structural reasons, but a gated SEIR head
stabilises every encoder. Code, pre-registrations, the experiment log and a
reproduction notebook accompany this paper.
